\documentclass[lang=english]{secseminar}
%\documentclass[lang=ngerman]{secseminar}

\usepackage{todonotes}

\title{Elementary Representation Theory}
\author{Fabian Weik}
\course{cs} % cs, se, mcl, msv, bis, simtech, or {my own special one}
\coursename{Hauptseminar Mathematik}
\coursesemester{WiSe 2022}
% Wenn der akademische Grad mit angegeben wird: Bitte auf korrekte Position achten (angelsächsische Grade nach dem Namen, deutsche Grade vor dem Namen).
\supervisor{Dr. Pascal Reisert}

\begin{document}

\begin{abstract}
    This document is an introduction to representation theory.
    The document starts with an introduction to the idea and applications of representation theory.
    After that, the document introduces all lesser-known mathematical concepts needed to understand representation theory, as well as facts used in later proofs.
    Then the basics of representation theory are presented.
    Such as the definition of representations, their properties, and equivalence of representations.
    After that, we will take a look at how representations can be classified as this brings structure into the set of representations a group can have.
    Here, the most important results of representation theory are presented.
    At last, we will classify the representations of the symmetric group $S_3$ as an example.
    This will also give us an application for cryptography.
\end{abstract}


\section{Introduction}

Representation theory is a branch of mathematics, that investigates the structure of groups.
It does so by \textit{representing} group elements as linear maps between vector spaces.
Since group elements are thought of as actions acting on something, it makes sense to represent them by maps acting on elements of vector spaces.
The idea is to represent group elements as actions on some vector space, such that the composition of these actions reflects the composition of the according group elements.
The vector space is then also called a $G$-module, and the mapping of the group elements to the linear maps is called the representation.

Linear algebra is a well-studied branch of mathematics.
And there are many algorithms for tackling problems in this field numerically.
Representing groups as linear maps, therefore, makes it easier to study these groups.
For example, Burnside's theorem was proven using representation theory in 1904~\cite{burnside1904groups}.
It took roughly 70 years to find a proof without using representation theory.

There are representations of groups using infinite-dimensional vector spaces, as well as finite-dimensional vector spaces.
In this document, we will focus on the case of finite-dimensional vector spaces since this is easier to understand and most readers are already familiar with finite-dimensional vector spaces.
The vector spaces used in this document are of the form $\C^n$ for reasons discussed later.

Representation theory finds applications anywhere groups are ubiquitous, i.e. mathematics, physics, chemistry etc.
The theory is also applicable in cryptography where it is used to study groups used in established protocols, such as the family of elliptic curve protocols.
As well as exploring groups to be used in emerging protocols.
An example of such emerging protocols are non-commutative group-based protocols.
At the end of this document, we will see why it is unwise to use the symmetric groups $S_n$ with some key-exchange protocols~\cite{khovanov2022monoidal}.

\section{Related Work}

This document is an introductory document for the topic of representation theory.
Therefore, the majority of works used are also introductory documents.

Fulton and Harris write a book called ``Representation theory: a first course''~\cite{fulton2013}.
It is a good introduction to the topic of representation theory and lays out the basics.
It does only cover finite-dimensional representations.

A similar work in german language is the book ``Einführung in die Struktur- und Darstellungstheorie der klassischen Gruppen'' by Hein~\cite{hein2013}.
This book also covers definitions of groups and vector spaces in depth.
So it is better for people that don't have a lot of knowledge of vector spaces.
The following section will also cover the fundamentals of group theory and linear algebra to remind the reader of these concepts.

\section{Fundamentals}

This section introduces various concepts, that we will rely upon in later sections.
Since representation theory is about representing group elements as linear maps vector spaces, the introduced concepts are from these two areas of mathematics.
The first subsection focuses on group theory.
The second subsection focuses on linear algebra.

\subsection{Group Theory}

A \textbf{group} $(G, \cdot)$ is a set $G$ with an operation $\cdot: G \times G \to G$ that satisfies the following axioms.
\begin{subequations}
\begin{align}
    \forall a, b, c \in G: \qquad & a \cdot (b \cdot c) = (a \cdot b) \cdot c & \text{associativity} \\
    \exists e \in G \forall g \in G: \qquad & e \cdot g = g \cdot e = g & \text{neutral element} \\
    \forall g \in G \exists g^{-1} \in G: \qquad & g^{-1} \cdot g = g \cdot g^{-1} = e & \text{inverse element}
\end{align}
If a group also satisfies commutativity, it is called an abelian group.
\begin{align}
    \forall a, b \in G: \qquad & a \cdot b = b \cdot a & \text{commutativity}
\end{align}
\end{subequations}

The \textbf{order} of a group element $a \in G$ is the number $n$, for which $a^n = 1$.
All elements in a finite group are of finite order.

A group can have multiple \textbf{conjugacy classes}.
Such a class is the set of all group elements that can be transformed into one another by conjugation.
The conjugation of $a \in G$ by $g \in G$ is defined as $g \cdot a \cdot g^{-1}$.
And therefore the conjugacy class of a group element $a \in G$ is
\begin{align}
    \Cl(a) & = \left\{ g \cdot a \cdot g^{-1} | g \in G \right\}
\end{align}

%\todo[inline]{remove topological group?}

%A \textbf{topological group} $(G, \cdot)$ is a group where both the operation $\cdot: G \times G \to G: (x, y) \mapsto xy$ and the inverse map $^{-1}: G \to G: x \mapsto x^{-1}$ are continuous.
%This is the case, if and only if the map $m: G \times G \to G: (x, y) \mapsto xy^{-1}$ is continuous.
%$m$ is continuous, if and only if for all $x, y \in G$ and any neighbourhood $W$ of $xy^{-1}$ there are neighbourhoods $U$ of $x$ and $V$ of $y$, such that $U \cdot V^{-1} \subseteq W$.

A \textbf{group homomorphism} is a map $\phi: G \to H$ between two groups $(G, \cdot)$ and $(H, *)$, such that
\begin{align}
    \phi(a \cdot b) = \phi(a) * \phi(b)
\end{align}
Or the diagram in \Cref{fig:fundamentals.groups.hom-cd} commutes.
Such diagrams are common in category theory and representation theory.
Understanding them will help in later sections.
\begin{figure}[h]
    \centering
    \begin{tikzcd}
        G \times G
            \arrow[r, "\phi \times \phi"]
            \arrow[d, "\cdot"]
        & H \times H
            \arrow[d, "*"]
        \\
        G   \arrow[r, "\phi"]
        & H
    \end{tikzcd}
    \caption{Commuting diagram for group homomorphisms}
    \label{fig:fundamentals.groups.hom-cd}
\end{figure}

A \textbf{field} $(F, +, -, 0, 1)$ is a set $F$ with two operations, addition $+: F \times F \to F$ and multiplication $\cdot: F \times F \to F$, such that $(F, +, 0)$ is an abelian group, $(F \setminus \{0\}, \cdot, 1)$ is a group, and the following axioms holds.
\begin{subequations}
\begin{align}
    \forall a, b, c \in F: \qquad & a \cdot (b + c) = a \cdot b + a \cdot c & \text{left distributive property} \\
    \forall a, b, c \in F: \qquad & (a + b) \cdot c = a \cdot c + b \cdot c & \text{right distributive property}
\end{align}
\end{subequations}
For this document, the only necessary field is the field of the complex numbers $(\C, +, \cdot, 0, 1)$.
When referring to this field in the following, only $\C$ is written.

\subsection{Linear Algebra}

A \textbf{vector space} over a field $F$ is a set $V$ with two operations, addition $+: V \times V \to V$ and scalar multiplication $\cdot: F \times V \to V$, that satisfy the following axioms.
\begin{align}
    \forall u, v, w \in V: \qquad & u + (v + w) = (u + v) + w & \text{associativity of vector addition} \\
    \forall u, v \in V: \qquad & u + v = v + u & \text{commutativity of vector addition} \\
    \exists 0_V \in V\ \forall v \in V: \qquad & v + 0_V = v & \text{identity element of vector addition} \\
    \forall v \in V\ \exists -v \in V: \qquad & -v + v = 0_V & \text{inverse element of vector addition} \\
    \forall v \in V\ a, b \in F: \qquad & a \cdot (b \cdot v) = (a \cdot b) \cdot v & \text{compatibility of multiplications} \\
    \forall v \in V: \qquad & 1 \cdot v = v & \text{identity element of scalar multiplication} \\
    \forall u, v \in V\ a \in F: \qquad & a \cdot (u + v) = a \cdot u + a \cdot v & \text{distributivity of scalar multiplication} \\
    \forall v \in V\ a, b \in F: \qquad & (a + b) \cdot v = a \cdot v + b \cdot v &  \text{distributivity of scalar multiplication}
\end{align}

A \textbf{linear map} $m$ between vector spaces $V$ and $W$ is a map $m: V \to W$ that satisfies the following two conditions.
\begin{align}
    \forall u, v \in V: \qquad & m(u + v) = m(u) + m(v) & \text{additivity} \\
    \forall c \in \C, v \in V: \qquad & m(c \cdot v) = c \cdot m(v) & \text{homogeneity}
\end{align}
In our case both vector spaces $V$ and $W$ always have the same, so $V = W$.
The set of all linear maps between $V$ and itself is denoted as $\GL(V)$.
\\

In most cases, the representations of groups are over finite-dimensional vector spaces.
If they are finite-dimensional, they will be vector spaces over the field of complex numbers $\C$ in this work.
A $n$-dimensional complex vector space is denoted as $\C^n$.
Its elements are written as vectors.

\todo[inline]{Matrices}
???
Such a map $M$ can be expressed via a \textbf{matrix}.
This is a core concept in representation theory since the elements of groups will be represented by matrices.



\newpage
\section{What are Representations}

This chapter introduces the concept of representations and some basic properties.
In the beginning, an abstract definition of representation is given, followed by a basic example.
After that, the notion of maps between representations is introduced, which is needed to define the equivalence of representations.

\subsection{Representations}

The \textbf{representation} $\rho$ of a group $G$ is defined as a group homomorphism $\rho: G \to \GL(V)$.
This means that $\rho$ preserves the structure of $G$.
So the diagram in \Cref{fig:main.what.rep-cd} commutes.
As you can infer from the diagram, the representation $\rho$ does not have to preserve \textit{all} structure of $G$.
The simplest representation therefore only has to preserve the structure, that every $g \in G$ is an element of $G$.
This is ridiculously trivial and you can see that it is a valid representation for any group $G$.
It is called the trivial representation and is defined as $\triv: G \to \GL(V): g \mapsto \id_V$.
\begin{figure}[h]
    \centering
    \begin{tikzcd}
        G \times G
            \arrow[r, "\rho \times \rho"]
            \arrow[d, "\cdot"]
        & \GL(V) \times \GL(V)
            \arrow[d, "\cdot"]
        \\
        G   \arrow[r, "\rho"]
        & \GL(V)
    \end{tikzcd}
    \caption{Commuting diagram for representations}
    \label{fig:main.what.rep-cd}
\end{figure}



\vspace{1cm}
\todo[inline]{Add main part. Might be multiple chapters}

\section{Conclusion}

\todo[inline]{Write Conclusion}


\vspace{1cm}
\todo[inline]{Add Outlook? Combined with Conclusion?}


\begin{abstract}
Im Seminar sammeln Sie erste Erfahrungen mit dem Erstellen von wissenschaftlichen Arbeiten.
Wir geben daher einen kurzen Überblick über die Ziele und Ihre Aufgaben beim Anfertigen einer
Seminararbeit.
Anschließend zeigen wir ein Beispiel für eine sinnvolle Struktur der Ausarbeitung.
Wir listen auf, welche Formalia Sie beachten sollten und geben zum Schluss einige wichtige Hinweise,
die Sie beachten sollten.
\end{abstract}

\section{Einleitung}
Ein Ziel des Studiums ist es, Studierenden das wissenschaftliche Arbeiten sowie
das wissenschaftliche Schreiben näher zu bringen. Erste Erfahrungen mit dem
Erstellen von wissenschaftlichen Arbeiten sammeln Studierende in der Regel
beim Verfassen von Seminarausarbeitungen und Abschlussarbeiten. Daher hier ein
kurzer Überblick, auf was beim Erstellen einer Seminarausarbeitung zu achten
ist.

\medskip Diese Kurzeinführung zum Thema Seminarausarbeitung basiert auf \cite{LIB,BOH07,KAR}.


\section{Ziele einer Seminararbeit}
Im Rahmen des Seminars ``Informationssicherheit und Kryptographie'' haben Sie die Aufgabe,
``Expert*in'' zu einem wissenschaftlichem Thema\footnote{Beachten Sie, dass Ihr Thema von Ihrem
Betreuer genauer eingegrenzt und Schwerpunkte definiert werden können.} zu werden. Dies bedeutet
insbesondere, dass
\begin{itemize}
\item Sie Ihr Thema inhaltlich verstehen,
\item Sie sich zu Ihrem Thema eigene Gedanken machen, und
\item Sie sich auch über die vorgegeben Unterlagen hinaus über Ihr Thema informieren.
\end{itemize}


Das Ziel der Seminararbeit ist, den Lesenden dieses fachliche Thema näher zu bringen und zu
erläutern. Folgende Punkte sind dabei insbesondere zu beachten:

\begin{itemize}
\item Ihre Aufgabe besteht nicht ausschließlich in der reinen Wiedergabe eines Themas. Auch die
  Ausarbeitung sollte zeigen, dass Sie sich eigene Gedanken zu Ihrem Thema gemacht haben.
\item Die Zielgruppe Ihrer Ausarbeitung sind Kommiliton*innen mit ähnlichem Wissensstand wie Sie. Das
  heißt, Sie können erwarten, dass Lesende Einführungsveranstaltungen im Bereich Informatik
  (einschließlich der Veranstaltung \emph{Grundlagen der Informationssicherheit}) erfolgreich
  besucht haben.
\item Ihre Ausarbeitung sollte in sich geschlossen sein.  Lesende müssen von Ihnen in die Lage
  versetzt werden, Ihre Ausarbeitung ohne Zuhilfenahme weiterer Literatur zu verstehen. Daher:
  Führen Sie grundlegende Begriffe und Konzepte ein, falls Sie nicht erwarten können, dass Lesende
  mit ihnen vertraut sind. Sollten Sie sich bei einem Aspekt nicht sicher sein, ob Sie deren Kenntnis
  bei den Lesenden voraussetzen können, fragen Sie Ihre Betreuungsperson.
\item Üblicherweise können Sie viel mehr zu einem Thema schreiben, als Sie Platz zur Verfügung
  haben. Es ist Teil Ihrer Aufgabe, das Thema in dem vorgegebenen Rahmen möglichst gut
  darzustellen. Dazu müssen Sie ggf. an einigen Stellen abstrahieren.
\item Vermitteln Sie neben den technischen Fakten und Sachverhalten den Lesenden auch eine Intuition
  für Ihr Thema. Veranschaulichungen, Beispiele und Bilder können --~in Maßen verwendet~-- Lesende
  beim Verstehen Ihres Themas unterstützen.  Finden Sie ein gutes Mittelmaß zwischen oberflächlicher
  Betrachtung und technischen Details.
\item Englische Fachbegriffe sollten nur sehr vorsichtig (oder besser gar nicht) eingedeutscht
  werden. Im Zweifel lieber zu wenig als zu viel eindeutschen bzw. gleich die ganze Ausarbeitung in
  englischer Sprache verfassen.
\end{itemize}

\section{Struktur der Ausarbeitung}
Eine wissenschaftliche Ausarbeitung sollte sinnvoll strukturiert sein. Das heißt, es sollte einen
roten Faden geben, der Lesende sinnvoll durch das Thema führt. Stellen Sie sicher, dass den Lesenden
zu jedem Zeitpunkt der Sinn und Zweck des aktuellen Kapitels klar ist, etwa indem Sie vorher einen
kurzen Überblick geben.

Im Folgenden zeigen wir Ihnen beispielhaft eine übliche Struktur für wissenschaftliche
Ausarbeitungen. Dabei sind Struktur, Reihenfolge, als auch Benennung der Abschnitte \emph{nicht
  strikt festgelegt}. Überlegen und entscheiden Sie selbst, wie Sie Ihr Thema verständlich und
technisch korrekt darstellen können.

\subsection{Kurzfassung / Abstract}
In der Kurzfassung (englisch: Abstract) geben Sie eine möglichst kurze und prägnante Inhaltsangabe
Ihrer Arbeit.
Sie sollte alle wichtigen Schlüsselwörter enthalten und eine ähnliche Struktur wie die Arbeit selber
haben.
Als Faustregel sollte jeder Abschnitt (``section'') in ca. einem Satz zusammengefasst werden.
Eine typische Länge für die Kurzfassung sind 150 bis 200 Wörter.

\subsection{Einleitung}
Eine Einleitung führt Lesende in das Thema der Ausarbeitung bzw. einer wissenschaftlichen Arbeit
ein. Erklären Sie, welches Problem Sie in Ihrer Ausarbeitung behandeln und warum dieses Problem ein
relevantes Problem ist. Vielleicht setzten Sie ihr Thema in einen sinnvollen Kontext. Geben Sie
einen Über- und Ausblick über die Ausarbeitung. Welcher Abschnitt behandelt welches Thema?

\subsection{Grundlagen}
Führen Sie alle Begriffe, die Ihre Lesergruppe potentiell nicht kennt, ein. Wählen Sie dabei die
passende Detailtiefe: Müssen Sie einen Begriff formal definieren oder reicht es, den Lesenden eine
Intuition zu vermitteln? Beschreiben Sie die Grundlagen prägnant.

\subsection{Hauptteil}
Der Hauptteil der Ausarbeitung ist der eigentliche Kern Ihrer Arbeit, in dem Sie ihr Thema den Lesenden
erklären. Der Begriff ``Hauptteil'' bezeichnet hierbei ein Konzept und ist keinesfalls die
Überschrift für diesen Abschnitt. Normalerweise umfasst der Hauptteil ohnehin mehrere Kapitel.

\subsection{Verwandte Arbeiten und Wissenschaftlicher Hintergrund}

Eine Auflistung verwandter Arbeiten ist essentiell für jeden wissenschaftlichen Artikel.\\

\noindent
\emph{Zur Position.}
Bereits in der Einleitung sollte auf die Unzulänglichkeiten verwandter Arbeiten eingegangen werden,
um die eigenen Arbeit zu motivieren.
In kürzeren Ausarbeitung mit einer kleinen Anzahl ähnlicher veröffentlichter wissenschaftlicher
Arbeiten kann man die verwandten Arbeiten alternativ schon komplett in der Einleitung abhandeln.
In längeren Ausarbeitungen oder sobald der Umfang der verwandten Literatur einen eigenen Abschnitt
rechtfertigt, sollte jedoch ein eigener Abschnitt "Related Work/Verwandte Arbeiten" angelegt werden.
In jedem Fall sollte das Thema "Verwandte Arbeiten" für Lesende erkennbar und nachvollziehar
dargestellt werden.
Der eigene Abschnitt kann entweder vor den eigentlichen Beitrag gestellt werden (also der Abschnitt
nach der Einleitung), um die Verbesserungen der Arbeit gegenüber anderen Arbeiten zu betonen.
Alternativ kann dieser Abschnitt auch erst am Ende eines Papers vor Zusammenfassung und Ausblick
auftauchen.\\

\noindent
\emph{Zum Inhalt}.
In den Abschnitt "Related Work" gehört beispielsweise ein bekannter Überblicks-Artikel zum
Themengebiet bzw. ein standard-referenziertes Buch.
Damit könn(t)en sich interessierte Forschende einen eigenen Überblick zu verwandten Arbeiten in
diesem Themengebiet schaffen, sofern dieser noch nicht vorhanden ist.
Es folgt die Auflistung ähnlicher bzw. verwandter Arbeiten.
Diese Arbeiten sollten dabei nicht einfach beschrieben werden, sondern immer in Bezug zur eigenen
Arbeit gesetzt werden.
Idealerweise sollten die Kernverbesserungen der eigenen Arbeit im Verhältnis zu jeder verwandten
Arbeit hier herausgearbeitet werden.
Auf keinen Fall sollten verwandte Arbeiten unreflektiert aufgelistet werden.

\subsection{Zusammenfassung und Ausblick}
In der Zusammenfassung fassen Sie die Resultate Ihrer Ausarbeitung zusammen. Hinterfragen Sie die
Resultate, die Sie im Rahmen Ihrer Ausarbeitung dargestellt haben. Setzen Sie die Resultate in einen
größeren Kontext. Geben Sie eine (eigene) Wertung über die Resultate.

\subsection{Referenzen/Literatur} \label{subsec:referenzen}

Hier listen Sie die Literatur auf, auf die Sie im Rahmen Ihrer Arbeit verweisen. Dieser Abschnitt
wird vom Tool \hologo{BibTeX} automatisch entsprechend den Templateregeln generiert. Erfassen Sie
Ihre Referenzen in der Datei \texttt{bibtex.bib} aus der Vorlage. \hologo{BibTeX}-Einträge für
wissenschaftliche Veröffentlichungen können Sie z.~B. von \url{https://dblp.org} (siehe auch
Abbildung \ref{abb:dblp})\footnote{Sollte DBLP Ihnen zwei \hologo{BibTeX}-Einträge ausgeben, so
  sollten Sie beide Einträge in das Literaturverzeichnis kopieren.} oder auch von
\url{https://scholar.google.de} exportieren. Weiter können Sie für die Verwaltung der
\hologo{BibTeX}-Einträge auch Software wie z.~B.
\emph{JabRef}\footnote{\url{https://www.jabref.org/}} verwenden.

\begin{figure}
    \centering
    \includegraphics[width=\textwidth]{dblp.png}
    \caption{\hologo{BibTeX} Referenzen aus DBLP exportieren.}
    \label{abb:dblp}
\end{figure}


\section{Formalia}
In diesem Abschnitt werden kurz die wichtigsten Eckpunkte zur Beachtung der Formalia einer
Seminararbeit dargestellt. Viele Formalia, wie etwa Schriftgröße und Schriftart, sind bereits durch
dieses (verpflichtende) Template vorgegeben und werden hier nicht separat erwähnt.


\subsection{Deckblatt}
Sie sollten folgende Regeln für das Erstellen des Deckblattes beachten:
\begin{itemize}
\item Geben Sie an, ob es sich um eine Seminarausarbeitung (Bachelor) oder eine
  Hauptseminarausarbeitung (Master) handelt und hinterlegen Sie die korrekte Veranstaltung.
\item Geben Sie Ihren Studiengang an (siehe auch Tabelle~\ref{tab:studiengangauswahl}).
\item Geben Sie Ihren Namen \textbf{und Matrikelnummer} an. Sofern Sie bereits einen akademischen
  Grad erreicht haben, geben Sie diesen mit an.
\end{itemize}

\begin{table}
    \begin{tabular}{l|l}
        \textbf{Eingabewert} & \textbf{Studiengang}\\
        \hline
        cs & Informatik\\
        se & Softwaretechnik\\
        mcl & Computerlinguistik\\ 
        tk & Technische Kybernetik\\
        msv & Maschinelle Sprachverarbeitung\\
        bis & Wirtschaftsinformatik\\
        simtech & Simulation Technology\\
        Studiengang & ``Studiengang'' (nicht vordefinierter Studiengang)\\ 
    \end{tabular}
    \caption{Deckblatt: Studiengangauswahl im Template im Befehl \texttt{\textbackslash course}}
    \label{tab:studiengangauswahl}
\end{table}


\subsection{Sprache, Orthografie und Grammatik}
Verfassen Sie die Ausarbeitung in Deutsch oder Englisch\footnote{Zu Beginn der \hologo{TeX}-Datei
  müssen Sie die Sprache entsprechend einstellen.}. Sowohl Orthografie als auch Grammatik fließen in
die Bewertung ein und sollten daher fehlerfrei sein.  Sollten Sie die Ausarbeitung in Englisch
verfassen, halten Sie sich an die in der Informatik übliche amerikanische Rechtschreibung.  Für die
meisten \LaTeX-fähigen Editoren (z.~B. TeXStudio) gibt es Erweiterungen zur Überprüfung der
Rechtschreibung -- wir empfehlen Ihnen, diese zu verwenden.

In einer deutschen Ausarbeitung verwenden Sie in der Informatik üblicherweise auch englische
Fachbegriffe -- sofern es keine etablierten deutschen Fachausdrücke als Gegenstück
gibt. Beispielsweise ist der Begriff ``Reihung'' den meisten Informatiker*innen nicht geläufig, während
jeder den entsprechenden englischen Begriff ``Array'' kennt.

\subsection{Zitieren}
Verweisen Sie im Rahmen Ihrer Ausarbeitung in den jeweiligen Textpassagen auf die Quellen, die Sie
verwendet haben bzw. deren Inhalt Sie darstellen.  Dabei sollten Ihre Literaturangaben das Auffinden
der Quellen möglichst einfach machen.  Beispiele für Zitierweisen:
\begin{itemize}
\item ... diese Resultate konnten bestätigt werden \cite{stickel2009wissenschaftliches}.
\item Wie Bohlinger \cite{BOH07} gezeigt hat, ...
\item Wie in \cite{KAR} dargestellt, ...
\item Der im folgenden besprochene Angriff wurde in \cite{DBLP:conf/ccs/AdrianBDGGHHSTV15} erstmalig
  beschrieben.
\end{itemize}
Vermeiden Sie durch geeignete Zitierweise, dass Sie in längeren Abschnitten jeden einzelnen Satz mit
der gleichen Referenz versehen.

\subsection{Tabellen, Grafiken und Code Fragmente}
Tabellen, Grafiken und Codefragmente sollten Sie jeweils mit einer Beschriftung versehen. Verweisen
Sie innerhalb Ihres Textes an der passenden Stelle auf das entsprechende Objekt. \hologo{LaTeX}
unterstützt Sie bei der korrekten Darstellung mit dem \texttt{table}, \texttt{graphicx} und
\texttt{listings}-Paket. Der Abschnitt~\ref{subsec:referenzen} gibt Ihnen ein Beispiel für die
Verwendung einer Grafik in der Ausarbeitung. (In wissenschaftlicher Literatur ist es üblich,
Grafiken nicht als Teil des Fließtextes einzubinden, sondern zu Beginn einer Seite oder auf einer
separaten Seite darzustellen -- die Platzierung übernimmt \hologo{LaTeX} für Sie.)

\section{Weitere Hinweise}
\begin{itemize}
\item Erstellen Sie Ihre Ausarbeitung \emph{vor} Ihrem Seminarvortrag und nutzen Sie die Zeit nach
  dem Vortrag zur Überarbeitung/Verbesserung der Ausarbeitung.
\item Bei Fragen und Unklarheiten, wenden Sie sich selbstständig an Ihre Betreuungsperson.
\item Ein guter Ausgangspunkt zum Finden von weiterführender Literatur sind unter anderem die
  Referenzen in ihrer Themenvorlage.
\item Stellen Sie sicher, dass Sie Konzepte einführen, bevor Sie diese verwenden. Dieser Hinweis
  gilt insbesondere auch für die Verwendung von Abkürzungen.
\item Nutzen Sie den Platz in der Ausarbeitung für die Darstellung von Inhalten sinnvoll
  aus. Vermeiden Sie Platzverschwendung (z.~B. durch schlecht skalierte Grafiken).
\item Überschriften sollten immer von Fließtext gefolgt sein, nicht unmittelbar von einer weiteren
  Unterüberschrift.
\item Eine Überschriftenebene sollten Sie nur mit Unterüberschriften unterteilen, wenn Sie
  mindestens zwei Unterabschnitte anlegen.
\item Fügen Sie \emph{keinen} manuellen Seitenumbruch nach Abschluss eines Abschnittes ein.
\item Hilfreich ist es, Ihre Arbeit Dritten, z.~B. Kommiliton*innen, zum Lesen zu geben. Sie können So
  wertvolles Feedback zum Inhalt und der Rechtschreibung Ihrer Ausarbeitung erhalten.
\end{itemize}


\bibliographystyle{plainnat}
\bibliography{bibtex}
\end{document}






%%% Local Variables:
%%% mode: latex
%%% TeX-master: t
%%% End:
